\documentclass[11pt]{article}
\usepackage[margin=1in]{geometry}
\usepackage{enumitem}
% =============================================================================
% Preambles/header.tex — shared LaTeX/Beamer preamble for this project
%
% Use from a Beamer lecture by prepending:
%     \input{header}
% and compiling with TEXINPUTS=../Preambles:$TEXINPUTS (the /compile-latex
% skill does this automatically).
%
% PALETTE CONTRACT. Color names defined here MUST match the names in
% Quarto/theme-template.scss so Beamer and Quarto renderings use the same
% palette. The scripts/check-palette-sync.sh script enforces this.
%
% To customize: edit the HEX values below AND the matching $variable in
% Quarto/theme-template.scss, then run scripts/check-palette-sync.sh.
% =============================================================================

% -----------------------------------------------------------------------------
% Engine detection + encoding
%
% Our /compile-latex skill uses XeLaTeX, where inputenc is unnecessary and
% can produce encoding warnings. Only load inputenc under pdfTeX.
% -----------------------------------------------------------------------------
\usepackage{iftex}
\ifPDFTeX
  \usepackage[utf8]{inputenc}
\fi

\usepackage{amsmath, amssymb, amsthm}
\usepackage{booktabs}
\usepackage{graphicx}

% -----------------------------------------------------------------------------
% PALETTE — mirrors Quarto/theme-template.scss variable names.
% Load xcolor and define colors BEFORE anything that references them
% (notably hyperref, hypersetup, and the Beamer color assignments).
% -----------------------------------------------------------------------------
\usepackage{xcolor}

\definecolor{primary-blue}{HTML}{012169}      % headings, accents
\definecolor{primary-gold}{HTML}{B9975B}      % emphasis, borders
\definecolor{highlight-yellow}{HTML}{F2A900}  % markers, alerts
\definecolor{light-bg}{HTML}{E8EDF5}          % title / section backgrounds
\definecolor{jet}{HTML}{1A1A1A}               % body text

% Semantic colors (used in diagrams and annotations)
\definecolor{positive}{HTML}{15803D}          % good / observed / identified
\definecolor{negative}{HTML}{B91C1C}          % bad / problematic / bias
\definecolor{neutral}{HTML}{525252}           % reference / context

% Highlight accents (match .hi-* classes in the SCSS)
\definecolor{hi-slate}{HTML}{314F4F}
\definecolor{hi-green}{HTML}{2E8B57}
\definecolor{hi-red}{HTML}{C41E3A}

% Semantic aliases so skill templates and snippets can write
% \textcolor{accent}{...} without hard-coding "primary-blue".
\colorlet{accent}{primary-blue}
\colorlet{accent2}{primary-gold}

% -----------------------------------------------------------------------------
% Hyperref — loaded AFTER xcolor + palette so link colors resolve.
% -----------------------------------------------------------------------------
\usepackage{hyperref}
\hypersetup{colorlinks=true, linkcolor=primary-blue, urlcolor=primary-blue,
            citecolor=primary-blue, pdfborder={0 0 0}}

% -----------------------------------------------------------------------------
% Course code — set once per deck (after \input{header}, before \begin{document})
% via \coursecode{CS401}. Shown in the footer on every slide so a deck's course
% is identifiable without opening the title page. Empty by default.
% -----------------------------------------------------------------------------
\makeatletter
\newcommand{\coursecode}[1]{\gdef\@coursecodetext{#1}}
\gdef\@coursecodetext{}
\makeatother

% -----------------------------------------------------------------------------
% Beamer theme assignments (only apply under Beamer).
%
% We detect Beamer via \@ifclassloaded{beamer}, which is the documented,
% non-internal way. This must be wrapped in \makeatletter/\makeatother
% because \@ifclassloaded contains an @-catcode character.
% -----------------------------------------------------------------------------
\makeatletter
\@ifclassloaded{beamer}{%
  \usefonttheme{professionalfonts}
  \setbeamercolor{structure}{fg=primary-blue}
  \setbeamercolor{frametitle}{fg=primary-blue}
  \setbeamercolor{title}{fg=primary-blue}
  \setbeamercolor{subtitle}{fg=primary-gold}
  \setbeamercolor{author}{fg=primary-blue}
  \setbeamercolor{institute}{fg=neutral}
  \setbeamercolor{date}{fg=primary-gold}
  \setbeamercolor{itemize item}{fg=primary-blue}
  \setbeamercolor{itemize subitem}{fg=primary-gold}
  \setbeamercolor{alerted text}{fg=primary-gold}
  \setbeamercolor{block title}{fg=white, bg=primary-blue}
  \setbeamercolor{block body}{bg=light-bg}
  % Semantic block variants: block=definition/concept (blue), exampleblock=
  % worked example (green/positive), alertblock=key takeaway or caution (gold).
  % Keep to <=2 colored boxes per slide across ALL three variants (INV-7).
  \setbeamercolor{block title example}{fg=white, bg=positive}
  \setbeamercolor{block body example}{bg=light-bg}
  \setbeamercolor{block title alerted}{fg=white, bg=primary-gold}
  \setbeamercolor{block body alerted}{bg=light-bg}

  % Minimal footer: course code (left) + page X / N (right), no navigation clutter
  \setbeamertemplate{navigation symbols}{}
  \setbeamertemplate{footline}{%
    \color{neutral}\footnotesize\hspace{0.7em}\@coursecodetext\hfill
    \insertframenumber/\inserttotalframenumber\hspace{0.7em}\vspace{0.3em}}
}{}
\makeatother

% -----------------------------------------------------------------------------
% TikZ setup — libraries the snippet gallery uses
% -----------------------------------------------------------------------------
\usepackage{tikz}
\usetikzlibrary{arrows.meta, positioning, calc, decorations.pathreplacing,
                fit, shapes.geometric, backgrounds}

% Shared TikZ styles — snippets and hand-written diagrams can pick these up
% via `\begin{tikzpicture}[every node/.style={...}]` or by naming specific
% styles (e.g., `\node[dag-node] (X) at (0,0) {X};`).
\tikzset{
  % Node styles for causal DAGs and similar
  dag-node/.style = {
    draw=primary-blue, fill=light-bg, rounded corners=2pt,
    minimum width=1.2cm, minimum height=0.9cm, align=center, font=\small
  },
  decision-node/.style = {
    draw=primary-gold, fill=white, diamond, aspect=1.8,
    minimum width=1.8cm, minimum height=1.2cm, align=center, font=\small,
    inner sep=1pt
  },
  % Edge styles — observed vs counterfactual
  observed-edge/.style = {-{Latex[length=2.5mm]}, thick, draw=jet},
  counterfactual-edge/.style = {-{Latex[length=2.5mm]}, thick, draw=neutral, dashed},
  confound-edge/.style = {-{Latex[length=2.5mm]}, thick, draw=negative, dashed},
  % Data-point styles for plots
  observed-dot/.style = {fill=primary-blue, draw=primary-blue, circle, inner sep=1.2pt},
  counterfactual-dot/.style = {fill=white, draw=neutral, circle, inner sep=1.2pt}
}

% -----------------------------------------------------------------------------
% Convenience macros
% -----------------------------------------------------------------------------
\newcommand{\muted}[1]{\textcolor{neutral}{#1}}
\newcommand{\key}[1]{\textcolor{primary-gold}{\textbf{#1}}}
\newcommand{\good}[1]{\textcolor{positive}{#1}}
\newcommand{\bad}[1]{\textcolor{negative}{#1}}

% Standout frame macro: full-bleed dark background for section transitions.
% Beamer-only (uses \begin{frame}/\setbeamercolor/\usebeamerfont) -- do not
% call from an article-class file (e.g. Notes/<CODE>/*-notes.tex). \muted,
% \key, \good, \bad, and \coursecode above are plain xcolor/gdef and are
% safe to use from either class; verified when Notes/article-class support
% was added.
% Expand: \transitionslide{Your transition title}
\newcommand{\transitionslide}[1]{%
  \begingroup
  \setbeamercolor{background canvas}{bg=primary-blue}
  \begin{frame}[plain]
    \centering
    \vfill
    {\usebeamerfont{title}\color{white}\Large #1\par}
    \vfill
  \end{frame}
  \endgroup
}

% Section divider frame (Beamer-only), an alternative to \transitionslide with a
% darker two-line style: near-black (jet) background, a small white label line
% above a larger gold title, and a thin gold rule along the bottom edge.
% Expand: \sectiondivider{Section 2}{Pointers}
\newcommand{\sectiondivider}[2]{%
  \begingroup
  \setbeamercolor{background canvas}{bg=jet}
  \begin{frame}[plain]
    \vspace*{3.0cm}
    {\color{white}\ttfamily\large #1\par}
    \vspace{0.5cm}
    {\color{primary-gold}\LARGE #2\par}
    \begin{tikzpicture}[remember picture, overlay]
      \draw[primary-gold, line width=2.5pt]
        ($(current page.south west)+(0,0.1)$) -- ($(current page.south east)+(0,0.1)$);
    \end{tikzpicture}
  \end{frame}
  \endgroup
}

% -----------------------------------------------------------------------------
% Channel watermark — "Concepts That Click" (YouTube).
%
% Article-class files (CompetitiveExam Q/A sets, Notes, Assignments) get a
% light diagonal draft watermark on every page plus a centered footer naming
% the channel. Beamer decks get a subtle TikZ background watermark instead
% (the footline already carries the course code). Centralized here so every
% MCQ PDF is branded without per-file edits.
% -----------------------------------------------------------------------------
\makeatletter
\@ifclassloaded{article}{%
  \usepackage{draftwatermark}
  \SetWatermarkText{Concepts That Click}
  \SetWatermarkScale{1}
  \SetWatermarkFontSize{2cm}
  \SetWatermarkLightness{0.86}
  \SetWatermarkAngle{45}
  \usepackage{fancyhdr}
  \pagestyle{fancy}
  \fancyhf{}
  \fancyfoot[C]{\footnotesize\textcolor{neutral}{Concepts That Click~|~YouTube: \href{https://www.youtube.com/@concepts.thatclick}{@concepts.thatclick}}}
  \renewcommand{\headrulewidth}{0pt}
  \renewcommand{\footrulewidth}{0pt}
  \fancypagestyle{plain}{%
    \fancyhf{}%
    \fancyfoot[C]{\footnotesize\textcolor{neutral}{Concepts That Click~|~YouTube: \href{https://www.youtube.com/@concepts.thatclick}{@concepts.thatclick}}}%
    \renewcommand{\headrulewidth}{0pt}%
    \renewcommand{\footrulewidth}{0pt}%
  }
}{}
\@ifclassloaded{beamer}{%
  \setbeamertemplate{background}{%
    \begin{tikzpicture}[remember picture, overlay]
      \node[rotate=45, opacity=0.055, align=center]
        at (current page.center)
        {\Huge\textcolor{neutral}{Concepts That Click}};
    \end{tikzpicture}%
  }
}{}
\makeatother 
\coursecode{JKSSB}
\renewcommand{\thesection}{14.\arabic{section}}
\newtheorem{example}{Example}[section]
\newenvironment{definitionbox}[1]{%
  \par\noindent\textbf{Definition (#1).}\ \itshape
}{\par\vspace{0.3em}}
\title{Output Devices --- Detailed Lecture Notes}
\author{JKSSB Computer Awareness}
\date{\today}

\begin{document}
\maketitle

\section{What is an output device}
A computer works in a cycle: input, process, and output. In the output step,
the result of processing is handed back to the user or to another device. An
\textbf{output device} is the hardware that receives processed data from the
CPU and presents it in a form a human can see, hear, or read. Without an
output device the result of the computation would stay locked inside the
machine.

\begin{definitionbox}{Output device}
Hardware that receives processed data from the CPU and presents it to the
user or to another device in a readable form, such as on a screen, through a
speaker, or on paper.
\end{definitionbox}

Common output devices include the monitor, projector, speaker, headphones,
printer, and plotter. The following table fixes the meaning of each term used
in this lesson.

\begin{center}
\begin{tabular}{p{0.24\textwidth}p{0.66\textwidth}}
\toprule
\textbf{Term} & \textbf{Meaning in this lesson} \\
\midrule
Monitor & The main visual output device; shows text, images, and video as soft copy. \\
Projector & Throws a large image onto a screen or wall as soft copy. \\
Speaker & Converts an electrical signal into sound for a group of listeners. \\
Headphones & Produce sound for a single listener; worn on or over the ears. \\
Printer & Produces hard copy on paper; impact or non-impact. \\
Plotter & Draws large, precise line drawings such as engineering designs. \\
MFP & Multifunction Printer; combines printing, scanning, copying, and often faxing. \\
\bottomrule
\end{tabular}
\end{center}

\section{Classifying output devices}
Output devices are grouped by the form in which they present the result. This
is the classification JKSSB question-setters use most often, so learn one
example for each row.

\begin{center}
\begin{tabular}{p{0.26\textwidth}p{0.60\textwidth}}
\toprule
\textbf{Category} & \textbf{Typical devices} \\
\midrule
Visual / soft copy & Monitor, projector \\
Audio & Speaker, headphones \\
Printed / hard copy & Printer, plotter \\
\bottomrule
\end{tabular}
\end{center}

A single device may belong to more than one row if it produces more than one
kind of output, but for exam purposes use the primary role: a monitor is a
visual output device, a speaker is an audio output device, and a printer is a
printed (hard copy) output device.

\section{Soft copy vs hard copy}
Output comes in two forms. \textbf{Soft copy} is output that is displayed
electronically, such as an image on a monitor or a projected slide; it is not
printed and can usually still be edited. \textbf{Hard copy} is output that is
printed on a physical medium such as paper, produced by a printer or a
plotter; once printed it is fixed.

\begin{definitionbox}{Soft copy}
Output that is shown on a screen or projected electronically and is not
printed, such as the image on a monitor or projector.
\end{definitionbox}

\begin{definitionbox}{Hard copy}
Output that is printed on a physical medium such as paper, produced by a
printer or a plotter.
\end{definitionbox}

The distinction matters because a single item can test it in either
direction. If the stem mentions a screen, a display, or a projected image,
the answer is soft copy; if it mentions paper, printing, or a physical
printout, the answer is hard copy.

\begin{center}
\begin{tabular}{p{0.20\textwidth}p{0.34\textwidth}p{0.34\textwidth}}
\toprule
\textbf{Point} & \textbf{Soft copy} & \textbf{Hard copy} \\
\midrule
Form & Electronic, on a screen & Physical, on paper \\
Examples & Monitor, projector & Printer, plotter \\
Editable & Easily edited & Fixed once printed \\
\bottomrule
\end{tabular}
\end{center}

\subsection{Uses and limits of each form}
Soft copy and hard copy serve different needs, and an exam item may ask which
is better for a stated purpose. Soft copy is convenient because it can be
edited, searched, copied, and shared electronically, and it costs nothing to
display; its limit is that it needs power and a screen, and it is easy to lose
if it is not saved. Hard copy is convenient because it needs no power to read,
is permanent, and can be signed or filed as a legal record; its limits are
that it can be bulky, is fixed once printed, and consumes paper and ink. The
rule of thumb is: to read, edit, or send, use soft copy; to sign, file, or
archive, use hard copy.

\section{Monitors}
A \textbf{monitor} is the main visual output device of a computer. It displays
text, images, and video as soft copy. Screen size is measured diagonally in
inches, and picture sharpness is measured in pixels, which together give the
resolution of the display. A monitor shows output to the user; it does not
take input, with the single exception of a touchscreen, which senses touch in
addition to displaying.

\subsection{CRT displays}
A \textbf{CRT} display uses a \textbf{Cathode Ray Tube}. An electron beam
sweeps across a phosphor-coated screen inside a vacuum tube to form the image.
CRT monitors are bulky and deep, and they are now largely replaced by flat
panels. For exam purposes the key point is the full form: CRT = \textbf{Cathode
Ray Tube}, and it is the older, bulky technology.

\subsection{Flat-panel displays: LCD, LED, and OLED}
A \textbf{flat-panel display} is thin, light, and flat. The common
technologies are LCD, LED, and OLED.

\begin{center}
\begin{tabular}{p{0.14\textwidth}p{0.44\textwidth}p{0.24\textwidth}}
\toprule
\textbf{Type} & \textbf{What it means} & \textbf{Backlight} \\
\midrule
LCD & Liquid Crystal Display; cells of liquid crystal control light & Needs a backlight \\
LED & Light Emitting Diode display; an LCD panel lit by LEDs & LED backlight \\
OLED & Organic Light Emitting Diode; each pixel emits its own light & No backlight (self-emitting) \\
\bottomrule
\end{tabular}
\end{center}

The nearest-neighbour trap is LCD versus LED versus OLED. An \textbf{LCD}
needs a separate backlight to be visible. An \textbf{LED} display is really an
LCD panel whose backlight is made of LEDs. An \textbf{OLED} display needs no
backlight at all, because each pixel makes its own light.

\subsection{Display size, resolution, and refresh}
The size of a monitor is given as the diagonal length of the screen in inches.
The picture is made of tiny dots called \textbf{pixels}; the number of pixels
across and down the screen is the \textbf{resolution}, such as 1920 by 1080. A
higher resolution shows more detail in the same screen area. The
\textbf{refresh rate} is how many times per second the screen is redrawn,
measured in hertz (Hz); a higher refresh rate gives smoother motion. These are
the three numbers a display specification quotes, and they apply to monitors
and to laptop screens alike.

\section{Projector}
A \textbf{projector} throws a large image onto a screen or a wall. It is used
in classrooms, meetings, and cinemas to show the same computer output to many
people at once. The projected image is soft copy, because it is shown on a
surface and not printed. A projector is an output device; it does not capture
anything, so it must not be confused with a scanner or a camera.

\section{Audio output: speakers and headphones}
A \textbf{speaker} converts an electrical signal into sound that a group of
listeners can hear. \textbf{Headphones} produce sound for a single listener
and are worn on or over the ears. Both are audio output devices.

The nearest-neighbour pair is the speaker against the microphone. A
\textbf{microphone} takes sound in and is an audio input device; a
\textbf{speaker} or a \textbf{headphone} sends sound out and is an audio
output device.

\begin{example}
A student plays a recorded lecture on the office desktop. The sound leaves the
computer through the speakers, so the speakers are audio output devices. If
the student then records a reply into the same computer, the microphone is
used; that is an audio input device. The same machine performs both jobs
through two different devices.
\end{example}

\section{Printers: overview}
A \textbf{printer} produces hard copy on paper. Printers fall into two broad
classes according to how they mark the paper. An \textbf{impact printer}
strikes a ribbon against the paper; a \textbf{non-impact printer} marks the
paper without striking it, using ink, toner, or heat. The common types are the
dot-matrix printer (impact) and the inkjet and laser printers (non-impact).
Laser printers are common in offices, while inkjet printers are common in
homes. A printer is an output device (TRAP-12).

\begin{center}
\begin{tabular}{p{0.18\textwidth}p{0.34\textwidth}p{0.34\textwidth}}
\toprule
\textbf{Point} & \textbf{Impact} & \textbf{Non-impact} \\
\midrule
Method & Strikes a ribbon and the paper & No striking; uses ink, toner, or heat \\
Examples & Dot-matrix & Inkjet, laser \\
Noise & Noisier & Quieter \\
Graphics quality & Lower & Higher \\
\bottomrule
\end{tabular}
\end{center}

A \textbf{dot-matrix} printer is an impact printer and is the noisiest of the
common types. An \textbf{inkjet} printer sprays tiny drops of ink onto the
paper, and a \textbf{laser} printer forms the image with toner that is fused
onto the paper by heat. Both are non-impact printers.

The choice follows the job. A dot-matrix printer is cheap to run and can print
multi-part carbon forms, but its quality is low and it is noisy. An inkjet
printer is quiet and inexpensive to buy and is suited to colour photographs
and home use, though the ink can be costly. A laser printer is fast, quiet,
and gives sharp text, which is why it is the usual office choice. Whatever the
type, every printer is an output device producing hard copy.

\begin{example}
An office prints a hundred letters every morning. The best choice is a laser
printer, because it is fast, quiet, and gives sharp text. If the same office
had to print only the occasional colour photograph at home, an inkjet printer
would be enough. If it needed multi-part carbon forms for a carbon copy, a
dot-matrix printer would be the right impact printer.
\end{example}

\section{Plotter}
A \textbf{plotter} draws large, precise line drawings, such as engineering and
architectural designs, on paper larger than a normal printer's. It draws using
a pen or a vector-based method, which makes it suitable for the continuous
lines of technical drawings. A plotter is an output device: it produces hard
copy but it does not read anything, so it must not be confused with a scanner
(TRAP-12).

\begin{example}
An item asks which output device is used to produce large engineering
drawings. The answer is a \textbf{plotter}, because it is designed for large,
precise line drawings. A normal printer is meant for ordinary pages, a monitor
shows soft copy, and a scanner is an input device.
\end{example}

\section{Multifunction printer (MFP)}
An \textbf{MFP} is a \textbf{Multifunction Printer}, also called an
\textbf{all-in-one} device. It combines \textbf{printing, scanning, copying},
and often \textbf{faxing} in a single unit. Because it does several jobs, an
MFP is the one common device that is both an input and an output device: its
scanner part takes input, while its printer part gives output. MFPs are common
in offices and small businesses.

\section{Summary of output devices}
The table below gathers the devices in this lesson for quick revision. Use it
as a checklist: for each device, be ready to state its category (visual,
audio, or printed) and its one-line function, and to name the device it is
most often confused with.

\begin{center}
\begin{tabular}{p{0.20\textwidth}p{0.22\textwidth}p{0.46\textwidth}}
\toprule
\textbf{Device} & \textbf{Category} & \textbf{Function} \\
\midrule
Monitor & Visual / soft copy & Shows text, images, and video on a screen. \\
Projector & Visual / soft copy & Throws a large image onto a screen or wall. \\
Speaker & Audio & Converts an electrical signal into sound for a group. \\
Headphones & Audio & Produces sound for a single listener. \\
Printer & Printed / hard copy & Prints text and graphics on paper; impact or non-impact. \\
Plotter & Printed / hard copy & Draws large, precise line drawings. \\
MFP & Hybrid & Combines printing, scanning, copying, and often faxing. \\
\bottomrule
\end{tabular}
\end{center}

\section{Points of confusion between similar devices}
Beyond classifying a device, JKSSB items often test a single distinction
between two devices that look or sound alike. The table below collects the
pairs used in this lesson.

\begin{center}
\begin{tabular}{p{0.30\textwidth}p{0.60\textwidth}}
\toprule
\textbf{Pair} & \textbf{The distinction to state} \\
\midrule
Printer vs plotter & Printer: standard pages of text and graphics. Plotter: large, precise line drawings. \\
Monitor vs projector & Monitor: image on its own screen. Projector: image thrown onto a screen or wall. \\
Speaker vs microphone & Speaker: audio output. Microphone: audio input. \\
Soft copy vs hard copy & Soft copy: shown on a screen. Hard copy: printed on paper. \\
LCD vs OLED & LCD: needs a backlight. OLED: self-emitting, no backlight. \\
Impact vs non-impact & Impact: strikes the paper (dot-matrix). Non-impact: no striking (inkjet, laser). \\
\bottomrule
\end{tabular}
\end{center}

\begin{example}
An office buys one machine that must print letters, scan signed forms, and
photocopy pages. The right choice is a Multifunction Printer (MFP), because it
does all three jobs in one unit. If the office instead needed very large
engineering plans, the correct device would be a plotter, and if it only
needed to show a presentation to a room, the correct device would be a
projector. Matching the output device to the job is the practical form of this
topic's classification.
\end{example}

\section{Other output devices on the exam}
A few output devices appear in real JKSSB items without being part of the
main list. A \textbf{sound card} produces audio output: it turns the digital
sound data into a signal for the speakers. A \textbf{Voice Response Unit
(VRU)} speaks responses to a caller and so gives audio output. \textbf{Haptic
gloves} and a \textbf{tactile display} give touch or force feedback to the
user and are therefore output devices. All of these feed information
\emph{back} to the user and must be told apart from input devices of similar
names, such as a \textbf{data glove}, which captures hand movement as input
(PYQ-09).

\begin{definitionbox}{Sound card}
An expansion card that produces audio output, converting digital sound data
into a signal that speakers or headphones can play.
\end{definitionbox}

\section{Exam retrieval and common confusions}
Seven distinctions prevent most errors on this topic.
\begin{enumerate}
  \item \textbf{Soft copy} is on a screen (monitor, projector); \textbf{hard
    copy} is printed on paper (printer, plotter). Do not swap them.
  \item A \textbf{plotter} and a \textbf{printer} are \textbf{output} devices;
    a \textbf{scanner}, a \textbf{light pen}, and a \textbf{digitizer} are
    \textbf{input} devices (TRAP-12).
  \item \textbf{CRT} = Cathode Ray Tube and is bulky; modern monitors are flat
    panels (LCD, LED, OLED).
  \item An \textbf{LCD} needs a backlight, an \textbf{LED} display is an LCD
    with an LED backlight, and an \textbf{OLED} display is self-emitting with
    no backlight.
  \item A \textbf{dot-matrix} printer is an \textbf{impact} printer; inkjet
    and laser printers are \textbf{non-impact}.
  \item A \textbf{plotter} draws large engineering drawings (PYQ-29), while a
    normal printer is meant for ordinary pages.
  \item A \textbf{sound card}, a \textbf{Voice Response Unit (VRU)}, a
    \textbf{haptic glove}, and a \textbf{tactile display} give output to the
    user; a \textbf{data glove} captures hand movement as input (PYQ-09).
\end{enumerate}

\begin{example}
An item lists four devices --- scanner, plotter, drum printer, and light pen
--- and asks which are output devices. Classify each one: the scanner is
input, the plotter is output, the drum printer is output, and the light pen is
input. The answer is the plotter and the drum printer. This is exactly the
input/output classification tested by TRAP-12.
\end{example}

\subsection*{Exercises}
\begin{enumerate}[label=\arabic*.]
  \item Define an output device and give three examples.
  \item Distinguish soft copy from hard copy, giving one example of each.
  \item Expand CRT, LCD, and OLED, and state which of them needs a backlight.
  \item Explain the difference between an impact printer and a non-impact
    printer, giving one example of each.
  \item Which output device is used for large engineering drawings, and why?
  \item State whether each is input or output: plotter, scanner, projector,
    digitizer tablet, speaker.
\end{enumerate}

\subsection*{Solutions}
\begingroup\small
\begin{enumerate}[label=\arabic*.]
  \item An output device receives processed data from the CPU and presents it
    to the user in a readable form. Examples: monitor, printer, speaker (any
    three valid devices).
  \item Soft copy is shown on a screen and is not printed, such as an image on
    a monitor or projector. Hard copy is printed on paper, such as a page from
    a printer or plotter.
  \item CRT = Cathode Ray Tube; LCD = Liquid Crystal Display; OLED = Organic
    Light Emitting Diode. The LCD needs a backlight; the OLED does not,
    because each pixel emits its own light.
  \item An impact printer strikes a ribbon against the paper, such as a
    dot-matrix printer. A non-impact printer marks the paper without striking
    it, such as an inkjet or a laser printer.
  \item A plotter is used, because it is designed to draw large, precise line
    drawings such as engineering and architectural designs (PYQ-29).
  \item Plotter --- output; scanner --- input; projector --- output; digitizer
    tablet --- input; speaker --- output.
\end{enumerate}
\endgroup
\end{document}
